\section{Extended Related Work}
\label{app:related-work}
\paragraph{Recovering transformations in existing representations.}
Lenc and Vedaldi \citeyearpar{lenc2015understanding} fit paired linear/affine maps with reconstruction and task losses and study cross-network stitching. Sun \citeyearpar{sun2026transformationlaws} studies compatible transformation families, encoder-fibre/kernel conditions, global/local realizations and error accumulation; its pooled linear-risk decomposition explicitly isolates a sharing cost. Its evaluations distinguish prediction, unseen-start transport, composition, and downstream fidelity in visual features. Action existence and shared-versus-local fitting therefore have direct precedents. Our contribution is the completed comparison of lexical/rank controls, coherent registrations, storage budgets and contextual transport in a fixed LLM panel; the supporting elementary identities are not a new general identification theory.

\paragraph{Calendar representations and their computational use.}
Linear concept geometry and its intervention metric are developed by \citet{park2023linear}, with categorical and hierarchical structure studied further by \citet{park2025categorical}. \citet{engels2025notall} identify multidimensional calendar features and causally patch circular representations in arithmetic tasks; their intervention-based appendix also discusses finite-group computations. Their statistical notion of feature irreducibility is distinct from invariant-subspace irreducibility in group representation theory. \citet{karkada2026emergent} connect language statistics to representational geometry. \citet{feucht2026arithmetic} investigate calendar arithmetic through a computational route involving ordinary addition and subsequent cyclic decoding. Consequently, our weekday experiment neither discovers calendar circles nor establishes a novel calendar-steering mechanism. It measures a particular fitted operator system on a fixed vocabulary and context distribution.

\paragraph{Group-operation mechanisms and competing interpretations.}
\citet{nanda2023progress} reverse engineer Fourier computations for modular addition. \citet{chughtai2023toy} recover representation-theoretic structure in networks trained on group products, while \citet{wu2025unified} develop compact performance proofs for a subset of group-operation networks and show that matching intervention effects need not identify an algorithm. Coset-based explanations \citep{stander2024cosets} and dihedral-multiplication mechanisms \citep{wei2026dihedral} supply further evidence that algebraic structure and noncommutative neural computations predate this work. Unlike those algorithmic tasks, our unchanged pretrained models are not trained by us to output binary group products. Multiplying recovered matrices externally and inserting one resulting activation edit is a different experiment from reverse engineering the unedited model's native multiplication algorithm.

\paragraph{Learning actions and measuring compositionality.}
Equivariance by Contrast \citep{schmidt2025equivariance} jointly learns latent coordinates and action representations from unlabeled action-paired observations. Its identifiability theorem assumes injective mixing, sufficiently diverse/spanning data and exact matching conditional distributions, and recovers coordinates up to one invertible linear map, with correspondingly conjugated actions. Here the encoder, readout and labeled input action are fixed; we assess predictive fits in those coordinates. This observation contract differs from joint latent recovery and does not extend its theorem or transfer its guarantees to our measurements.

Homomorphism autoencoders learn action representations from transitions \citep{keurti2023homomorphism}. \citet{an2026homomorphism} assess compositional generalization with representational homomorphism errors on expression tasks; this is a direct precedent for evaluating hidden representations through composition, rather than a geometry-only comparison. We fit element-indexed transport operators on contextual lexical states and separately test prediction and products. Equivariant architectures \citep{cohen2016group,kondor2018generalization,weiler2019general}, group crosscoders, and equivariant sparse autoencoders \citep{gorton2024group,erdogan2025group} provide related symmetry-aware tools, while making different architectural or dictionary-learning commitments.

\paragraph{Probes, equivalence, and causal evidence.}
Probe selectivity requires controls for the tasks a fitted map can learn \citep{hewitt2019designing}. \citet{yao2026groupaction} emphasize the equivalence relation underlying representation claims; the coordinate-change group in that discussion must not be confused with our semantic input-action group. Exact laws survive conjugacy, whereas the reported approximate Euclidean defects depend on the selected metric. \citet{sun2026tracking} compare implementation families through causal abstractions and Congruity, with explicit assumptions linking representation distances to interpretive equivalence. Our same-element intertwiner is a narrower compatibility fit, not that test. Causal abstraction requires intervention consistency beyond directional score movement, and successful control does not alone identify native reasoning \citep{geiger2022inducing,geiger2025causal}.

\paragraph{Geometry, steering, and fields.}
Word relations, spatial probes, and representational convergence provide broader context \citep{mikolov2013linguistic,bolukbasi2016man,gurnee2023language,huh2024platonic}. Activation and representation engineering \citep{turner2023activation,zou2023representation}, manifold steering \citep{wurgaft2026manifold}, and contextual transformation geometry \citep{hu2026language} motivate behaviorally meaningful transformations. Conditional-mean interaction steering also depends on a factorial construction \citep{arora2026causal}; a state-centroid control is only a restricted adaptation. Our constant global displacement estimator is not a reproduction of these methods. A linear operator itself admits the state-dependent residual form $zR_g=z+z(R_g-\Id)$; the tested distinction is constant displacement versus state-dependent transport, not additive versus algebraic structure in general.

